\subsection{60. Learning in Modal Space: Solving Stochastic PDEs with PINNs (Zhang et al.\ 2019)}
\textbf{arXiv:} \texttt{1905.01205}\quad\textbf{Authors:} Zhang, Lu, Guo, Karniadakis (2019)\quad\textbf{Domain:} ML / PDE / UQ\\
\scorebadge{Coverage}{7}\quad\scorebadge{Agreement}{5}\quad\textbf{Total:} 12/20\par\smallskip
\noindent\textbf{Coverage rationale.} All 3 examples attempted (stochastic advection, stochastic Burgers, reaction-diffusion). 10/13 testable claims evaluated (77\%). NN-DO and NN-BO methods reimplemented in PyTorch from paper description (no official code available). Analytical solutions independently verified for advection and Burgers. Eigenvalue crossings computed analytically ($\sim$30 crossings in $[0,10\pi]$).\par
\noindent\textbf{Agreement rationale.} Eigenvalue crossing handling (Claim~9) is verified---a key qualitative contribution. Variance approximations are within an order of magnitude (e.g., advection NN-DO: 1.14\% vs paper's 0.11\%). However, mean ($E[u]$) errors are 20--35$\times$ higher: advection NN-DO 44.98\% vs paper's 1.96\%; Burgers NN-DO 14.09\% vs paper's 0.40\%. Root causes: (1)~PINN training failed (scaling factors collapsed), requiring fallback to supervised training; (2)~modal decomposition gauge freedom; (3)~no official code to calibrate hyperparameters. Inverse problem (reaction-diffusion) not solved: PDE coefficients $a$, $b$ absent from supervised loss.\par\smallskip
\noindent\textbf{Key findings:}
\begin{itemize}[leftmargin=1.4em,itemsep=0.2em,topsep=0.2em]
\item The NN-DO/NN-BO mathematical framework is correct and produces qualitatively correct results
\item Eigenvalue crossings (key to NN-BO's advantage over NN-DO) independently confirmed
\item Exact error reproduction requires the paper's PINN training expertise and hyperparameter tuning, which is not described in sufficient detail
\item Supervised training gives a lower bound on achievable accuracy; our results suggest the paper's sub-1\% errors require very careful tuning
\end{itemize}
\noindent\textbf{Verdict: PARTIAL.} Mathematical framework sound; quantitative claims not reproduced.
\noindent\textbf{Follow-on research questions:}
\begin{enumerate}[leftmargin=1.6em,itemsep=0.2em,topsep=0.2em,label=Q\arabic*.]
\item Implement the full physics-informed loss (PDE residual + DO/BO constraints) with NTK-based loss weighting---does this recover sub-1\% errors for the advection example?
\item Can the gauge freedom in the modal decomposition be resolved by enforcing explicit orthonormality constraints during training? Compare constrained vs unconstrained optimization.
\item Extend NN-BO to higher stochastic dimensions (10--50 KL modes) and measure how accuracy scales---is there a curse of dimensionality?
\item Apply NN-DO/NN-BO to a stochastic PDE with non-Gaussian input uncertainty (e.g., log-normal permeability in Darcy flow)---do the claimed advantages persist?
\item Compare NN-DO/NN-BO against modern alternatives (DeepONet, neural operators) for stochastic PDEs on the same benchmark suite.
\end{enumerate}

